\documentclass[conference]{IEEEtran}
\usepackage{amsmath}
\title{KV Cache Migration Revisited: A Measurement Study}
\author{Author B}
\begin{document}
\maketitle

\begin{abstract}
This paper measures key-value cache migration across edge nodes. We collected one week of traces from a small cluster and report transfer sizes, link utilization, and recomputation costs. Our measurements suggest that migration is often bandwidth-bound, and that recomputation is cheaper than expected for short contexts. We release the trace summaries to support reproducibility.
\end{abstract}

\section{Introduction}
Key-value cache migration is widely assumed to be expensive. We test that assumption with measurements rather than models. Our dataset covers five nodes and three link types. The main finding is that context length dominates transfer cost, while link type dominates variance. We also observe that many migrations are triggered by policies that ignore recomputation cost.

\section{Measurement Setup}
We instrumented a five-node cluster over seven days. Each node logged cache size, transfer duration, link utilization, and recomputation events. We did not modify the serving engine. Traffic was generated by a closed-loop workload with fixed request rates.

\section{Results}
Transfer duration scales linearly with context length. Link utilization spikes during migration and decays within seconds. Recomputation costs less than migration for contexts below 2K tokens. For longer contexts the ranking reverses. We found no evidence that node heterogeneity matters at this scale.

\section{Discussion}
Our measurements contradict the common assumption that migration is always preferable to recomputation. Policies should condition on context length. The result is limited to a single cluster and a synthetic workload.

\section{Conclusion}
Migration should be treated as a bandwidth decision, not a default. We hope the trace summaries help others calibrate their policies.

\end{document}
